\subsection{Cascode Mirror}

\subsubsection{Overview}
Constructed from four transistors, It improves
output resistance and accuracy by reducing the
effects of channel length modulation.

\subsubsection{Schematic}
\begin{circuitfig}
  % Paths, nodes and wires:
  \node[shape=rectangle, draw={rgb,255:red,255;green,0;blue,0}, line width=1pt, dash pattern={on 4pt off 2pt on 1pt off 2pt on 1pt off 2pt}, minimum width=1.715cm, minimum height=2.195cm] at (6.125, 7.135){};
  \draw (3.23, 8.04) to[american voltage source, l_={$V_{cc}$}] (1.98, 8.04);
  \node[nmos, arrowmos, xscale=-1] (N1) at (3.5, 2.75){} node[anchor=east] at (N1.text){$M_0$};
  \node[nmos, arrowmos] (N2) at (6, 2.75){} node[anchor=west] at (N2.text){$M_1$};
  \draw (3.5, 7.54) to[european resistor, l={$G_4$}] (3.5, 6.04);
  \draw (6, 5.79) |- (6, 5.25);
  \draw (5.02, 2.75) -- (4.48, 2.75);
  \draw (3.5, 1.98) -| (3.5, 1.75) -| (6, 1.98);
  \draw (4, 1) to[american voltage source, l={$V_{ee}$}] (2.75, 1);
  \draw (4, 1) -| (4.75, 1.75);
  \node[ground, rotate=-90] at (1.98, 8.04){};
  \node[ground, rotate=-90] at (2.75, 1){};
  \draw[dash pattern={on 0.4pt off 0.8pt}, -stealth] (6, 5.5) -- (6.5, 5.5);
  \node[shape=rectangle, minimum width=0.965cm, minimum height=0.465cm] (N3) at (7, 5.46){} node[anchor=center] at (N3.text){$V_{out}$};
  \draw (6, 7.5) to[european resistor, l={\textcolor{rgb,255:red,255;green,0;blue,0}{$G_5$}}] (6, 6);
  \draw (6, 5.79) -- (6, 6.02);
  \draw (3.5, 7.54) |- (3.23, 8.04);
  \draw (3.5, 8.04) -| (6, 7.5);
  \draw (10.25, 10.75) -- (10.25, -0);
  \node[shape=rectangle, minimum width=5.59cm, minimum height=1.09cm] at (5.438, 9.937){} node[anchor=north, align=center, text width=5.202cm, inner sep=6pt] at (5.438, 10.5){Large Signal Schematic\\(Sink)};
  \node[shape=rectangle, minimum width=3.09cm, minimum height=2.465cm] at (8.687, 7){} node[anchor=north west, align=left, text width=2.702cm, inner sep=6pt] at (7.125, 8.25){\textcolor{rgb,255:red,255;green,0;blue,0}{\footnotesize Exists Only For Analysis. In  Implementation take output directly from V\_\{out\}}};
  \node[shape=rectangle, minimum width=0.965cm, minimum height=0.465cm] (N4) at (7, 3.71){} node[anchor=center] at (N4.text){$V_{y}$};
  \draw[dash pattern={on 0.4pt off 0.8pt}, -stealth] (6, 3.71) |- (6.5, 3.71);
  \node[shape=rectangle, minimum width=0.965cm, minimum height=0.465cm] (N5) at (4.75, 4.25){} node[anchor=center] at (N5.text){$I_{G_1}$};
  \draw (17.625, 2.75) to[american current source, l_={$g_{m_1}V_{gs}$}] (17.625, 1.25);
  \draw (19.125, 2.75) to[european resistor, l_={$g_{o_1}$}] (19.125, 1.25);
  \draw (17.625, 1.25) -- (19.125, 1.25);
  \draw (19.125, 2.75) -- (17.625, 2.75);
  \draw (17.625, 1.25) -- (16.375, 1.25);
  \draw (16.375, 2.75) -- (15.625, 2.75);
  \node[shape=rectangle, minimum width=0.465cm, minimum height=0.465cm] at (16.375, 2.5){} node[anchor=center, align=center, text width=0.077cm, inner sep=6pt] at (16.375, 2.5){\footnotesize G};
  \node[shape=rectangle, minimum width=0.465cm, minimum height=0.465cm] at (16.375, 1.5){} node[anchor=center, align=center, text width=0.077cm, inner sep=6pt] at (16.375, 1.5){\footnotesize S};
  \node[shape=rectangle, minimum width=0.465cm, minimum height=0.465cm] at (19.625, 2.5){} node[anchor=center, align=center, text width=0.077cm, inner sep=6pt] at (19.625, 2.5){\footnotesize D};
  \draw (15.625, 2.75) -- (14.875, 2.75);
  \draw (13.375, 2.75) to[american current source, l={$g_{m_0}V_{gs}$}] (13.375, 1.25);
  \draw (11.875, 2.75) to[european resistor, l={$g_{o_0}$}] (11.875, 1.25);
  \draw (13.375, 2.75) -- (11.875, 2.75);
  \draw (13.375, 1.25) -- (11.875, 1.25);
  \draw (11.875, 2.75) -| (11.875, 4.25);
  \draw (14.875, 1.25) -- (13.375, 1.25);
  \node[shape=rectangle, minimum width=0.465cm, minimum height=0.465cm] at (14.875, 2.5){} node[anchor=center, align=center, text width=0.077cm, inner sep=6pt] at (14.875, 2.5){\footnotesize G};
  \node[shape=rectangle, minimum width=0.465cm, minimum height=0.465cm] at (14.875, 1.5){} node[anchor=center, align=center, text width=0.077cm, inner sep=6pt] at (14.875, 1.5){\footnotesize S};
  \node[shape=rectangle, minimum width=0.465cm, minimum height=0.465cm] at (11.375, 2.5){} node[anchor=center, align=center, text width=0.077cm, inner sep=6pt] at (11.375, 2.5){\footnotesize D};
  \draw (11.875, 8.562) to[european resistor, l={$G_4$}] (11.875, 7.062);
  \draw (11.875, 7.062) -| (11.875, 6.562);
  \draw (11.875, 8.562) -| (11.875, 9.062) -- (12.625, 9.062) -| (12.625, 8.812);
  \draw (19.125, 7.75) to[european resistor, l={\textcolor{rgb,255:red,255;green,0;blue,0}{$G_5$}}] (19.125, 6.25);
  \draw (19.125, 6.25) -| (19.125, 5.75);
  \draw (19.125, 7.75) -| (19.125, 8.25) |- (18.375, 8.25) -| (18.375, 8);
  \node[shape=rectangle, minimum width=0.715cm, minimum height=0.465cm] (N6) at (20, 3.5){} node[anchor=center] at (N6.text){$V_{y}$};
  \node[shape=rectangle, minimum width=0.965cm, minimum height=0.465cm] (N7) at (20.125, 5.75){} node[anchor=center] at (N7.text){$V_{out}$};
  \node[ground] at (14.875, 1.25){};
  \node[ground] at (16.375, 1.25){};
  \node[ground] at (12.625, 8.812){};
  \node[ground] at (18.375, 8){};
  \node[shape=rectangle, minimum width=0.965cm, minimum height=0.465cm] (N8) at (16.625, 3.563){} node[anchor=center] at (N8.text){$I_{G_1}$};
  \node[shape=rectangle, minimum width=4.84cm, minimum height=1.09cm] at (15.25, 9.937){} node[anchor=north, align=center, text width=4.452cm, inner sep=6pt] at (15.25, 10.5){Small Signal Schematic\\(Sink)};
  \draw (6, 3.52) -| (6, 3.71);
  \node[nmos, arrowmos] (N9) at (6, 4.48){} node[anchor=west] at (N9.text){$M_2$};
  \node[shape=rectangle, minimum width=0.965cm, minimum height=0.465cm] (N10) at (2.5, 3.71){} node[anchor=center] at (N10.text){$V_{x}$};
  \draw[dash pattern={on 0.4pt off 0.8pt}, -stealth] (3.5, 3.71) |- (3, 3.71);
  \draw (17.625, 5.75) to[american current source, l_={$g_{m_2}V_{gs}$}] (17.625, 4.25);
  \draw (19.125, 5.75) to[european resistor, l_={$g_{o_2}$}] (19.125, 4.25);
  \draw (17.625, 4.25) -- (19.125, 4.25);
  \draw (19.125, 5.75) -- (17.625, 5.75);
  \draw (17.625, 4.25) -- (16.375, 4.25);
  \draw (16.375, 5.75) -- (15.625, 5.75);
  \node[shape=rectangle, minimum width=0.465cm, minimum height=0.465cm] at (16.375, 5.5){} node[anchor=center, align=center, text width=0.077cm, inner sep=6pt] at (16.375, 5.5){\footnotesize G};
  \node[shape=rectangle, minimum width=0.465cm, minimum height=0.465cm] at (16.375, 4.5){} node[anchor=center, align=center, text width=0.077cm, inner sep=6pt] at (16.375, 4.5){\footnotesize S};
  \draw (19.125, 4.25) -- (19.125, 2.75);
  \draw[dash pattern={on 0.4pt off 0.8pt}, -stealth] (19.125, 3.5) -- (19.625, 3.5);
  \node[shape=rectangle, minimum width=0.715cm, minimum height=0.465cm] (N11) at (11, 3.563){} node[anchor=center] at (N11.text){$V_{x}$};
  \draw[dash pattern={on 0.4pt off 0.8pt}, -stealth] (11.875, 3.563) -- (11.375, 3.563);
  \node[nmos, arrowmos, xscale=-1] (N12) at (3.5, 4.48){} node[anchor=east] at (N12.text){$M_3$};
  \draw (3.5, 3.71) -| (3.5, 3.52);
  \draw (4.48, 4.48) -- (5.02, 4.48);
  \draw (4.75, 4.5) |- (3.5, 5.5);
  \draw (3.5, 5.25) -| (3.5, 5.5);
  \draw (3.5, 6.04) -| (3.5, 5.5);
  \node[shape=rectangle, minimum width=0.965cm, minimum height=0.465cm] (N13) at (2.5, 5.5){} node[anchor=center] at (N13.text){$V_{z}$};
  \draw[dash pattern={on 0.4pt off 0.8pt}, -stealth] (3.5, 5.5) |- (3, 5.5);
  \node[shape=rectangle, minimum width=0.965cm, minimum height=0.465cm] (N14) at (4.75, 6){} node[anchor=center] at (N14.text){$I_{G_2}$};
  \draw (15.625, 5.75) -- (14.875, 5.75);
  \draw (13.375, 5.75) to[american current source, l={$g_{m_3}V_{gs}$}] (13.375, 4.25);
  \draw (11.875, 5.75) to[european resistor, l={$g_{o_3}$}] (11.875, 4.25);
  \draw (13.375, 5.75) -- (11.875, 5.75);
  \draw (13.375, 4.25) -- (11.875, 4.25);
  \draw (14.875, 4.25) -- (13.375, 4.25);
  \node[shape=rectangle, minimum width=0.465cm, minimum height=0.465cm] at (14.875, 5.5){} node[anchor=center, align=center, text width=0.077cm, inner sep=6pt] at (14.875, 5.5){\footnotesize G};
  \node[shape=rectangle, minimum width=0.465cm, minimum height=0.465cm] at (14.875, 4.5){} node[anchor=center, align=center, text width=0.077cm, inner sep=6pt] at (14.875, 4.5){\footnotesize S};
  \node[shape=rectangle, minimum width=0.465cm, minimum height=0.465cm] at (11.375, 5.5){} node[anchor=center, align=center, text width=0.077cm, inner sep=6pt] at (11.375, 5.5){\footnotesize D};
  \draw (15.625, 5.75) |- (11.875, 6.562) |- (11.875, 5.75);
  \node[shape=rectangle, minimum width=0.715cm, minimum height=0.465cm] (N15) at (11, 6.562){} node[anchor=center] at (N15.text){$V_{z}$};
  \draw[dash pattern={on 0.4pt off 0.8pt}, -stealth] (11.875, 6.562) -- (11.375, 6.562);
  \node[shape=rectangle, minimum width=0.965cm, minimum height=0.465cm] (N16) at (16.625, 6.562){} node[anchor=center] at (N16.text){$I_{G_2}$};
  \draw[dash pattern={on 0.4pt off 0.8pt}, -stealth] (19.125, 5.75) -- (19.625, 5.75);
  \draw[dash pattern={on 0.4pt off 0.8pt}, -stealth] (4.75, 5.5) -- (4.75, 5.75);
  \draw[dash pattern={on 0.4pt off 0.8pt}, -stealth] (4.75, 3.71) -- (4.75, 4);
  \draw[dash pattern={on 0.4pt off 0.8pt}, -stealth] (15.625, 6.562) -- (16.125, 6.562);
  \draw[dash pattern={on 0.4pt off 0.8pt}, -stealth] (15.625, 3.563) -- (16.125, 3.563);
  \node[shape=rectangle, draw={rgb,255:red,255;green,0;blue,0}, line width=1pt, dash pattern={on 4pt off 2pt on 1pt off 2pt on 1pt off 2pt}, minimum width=1.215cm, minimum height=2.195cm] at (19.375, 7.385){};
  \draw (3.5, 3.71) -| (4.75, 2.75);
  \draw (15.625, 2.75) |- (11.875, 3.563);
\end{circuitfig}

\subsubsection{Transfer Function}
Here's a table of the voltages at each terminal of the transistors:

\begin{tabularx}{\linewidth}{|X|X|X|X|X|X|} \hline
  Transistor & Gate (G) & Source (S) & Drain (D) & $V_{GS}$ & $V_{DS}$ \\ \hline
  $M_0$ & $V_{X}$ & $V_{ee}$ & $V_{X}$ & $V_{X} - V_{ee}$ & $V_{X} - V_{ee}$ \\ \hline
  $M_1$ & $V_{X}$ & $V_{ee}$ & $V_{Y}$ & $V_{X} - V_{ee}$ & $V_{Y} - V_{ee}$ \\ \hline
  $M_2$ & $V_{Z}$ & $V_{Y}$ & $V_{OUT}$ & $V_{Z} - V_{Y}$ & $V_{OUT} - V_{Y}$ \\ \hline
  $M_3$ & $V_{Z}$ & $V_{X}$ & $V_{Z}$ & $V_{Z} - V_{X}$ & $V_{Z} - V_{X}$ \\ \hline
\end{tabularx}

\begin{enumerate}
  \item \textbf{Large Signal, Time Domain, DC Input Analysis (Q5)}:\\
    Applying KCL at $V_{OUT}$:
    \begin{align*}
      (V_{OUT} - V_{cc})G_5 + I_{M_2} &= 0 \\
      V_{OUT}G_5 - V_{cc}G_5 + I_{M_2} &= 0 \\
      V_{OUT} = V_{cc} - \frac{I_{M_2}}{G_5} \tag{1}
    \end{align*}
    Applying KCL at $V_{X}$:
    \begin{align*}
      I_{M_0} + I_{G_1} - I_{M_3} &= 0 \\
      I_{M_0} + I_{G_1} = I_{M_3} \\
      I_{M_0} &\approx I_{M_3} \quad \text{[assuming } I_{G_1} \approx 0 \text{ by MOS physics]} \tag{2}
    \end{align*}
    Applying KCL at $V_{Y}$:
    \begin{align*}
      I_{M_1} - I_{M_2} &= 0 \\
      I_{M_1} = I_{M_2} \tag{3}
    \end{align*}
    Applying KCL at $V_{Z}$:
    \begin{align*}
      (V_{Z} - V_{cc})G_4 + I_{M_3} + I_{G_2} &= 0 \\
      V_{Z}G_4 - V_{cc}G_4 + I_{M_3} + I_{G_2} &= 0 \\
      V_{Z}G_4 &= V_{cc}G_4 - (I_{M_3} + I_{G_2}) \\
      V_{Z} &= V_{cc} - \frac{(I_{M_3} + I_{G_2})}{G_4} \tag{4}
    \end{align*}
    Using the transistor equations for $I_{M_0}$, $I_{M_1}$, $I_{M_2}$, and $I_{M_3}$:
    \begin{align*}
      I_{M_0} &= k_{n_0} (V_{GS_0} - V_{th})^2 (1 + K_{CLM} V_{DS_0}) \\
              &= k_{n_0} (V_{X} - V_{ee} - V_{th})^2 (1 + K_{CLM} (V_{X} - V_{ee})) \\
              &= k_{n_0} (V_{X} - V_{ee} - V_{th})^2 \quad \text{[assuming } K_{CLM} \approx 0 \text{]} \tag{5} \\
      I_{M_1} &= k_{n_1} (V_{GS_1} - V_{th})^2 (1 + K_{CLM} V_{DS_1}) \\
              &= k_{n_1} (V_{X} - V_{ee} - V_{th})^2 (1 + K_{CLM} (V_{Y} - V_{ee})) \\
              &= k_{n_1} (V_{X} - V_{ee} - V_{th})^2 \quad \text{[assuming } K_{CLM} \approx 0 \text{]} \tag{6} \\
      I_{M_2} &= k_{n_2} (V_{GS_2} - V_{th})^2 (1 + K_{CLM} V_{DS_2}) \\
              &= k_{n_2} (V_{Z} - V_{Y} - V_{th})^2 (1 + K_{CLM} (V_{OUT} - V_{Y})) \\
              &= k_{n_2} (V_{Z} - V_{Y} - V_{th})^2 \quad \text{[assuming } K_{CLM} \approx 0 \text{]} \tag{7} \\
      I_{M_3} &= k_{n_3} (V_{GS_3} - V_{th})^2 (1 + K_{CLM} V_{DS_3}) \\
              &= k_{n_3} (V_{Z} - V_{X} - V_{th})^2 (1 + K_{CLM} (V_{Z} - V_{X})) \\
              &= k_{n_3} (V_{Z} - V_{X} - V_{th})^2 \quad \text{[assuming } K_{CLM} \approx 0 \text{]} \tag{8}
    \end{align*}
    It can be seen that,
    \begin{align*}
      I_{IN} = I_{M_0} &= I_{M_3} \quad \text{[from (2)]} \\
      I_{OUT} = I_{M_1} &= I_{M_2} \quad \text{[from (3)]} \\
      I_{M_0} &= I_{M_1} \quad \text{[from (5) and (6) and assuming } k_{n_0} = k_{n_1} \text{]} \\
      \therefore, I_{IN} &= I_{OUT}
    \end{align*}
    Thus,
    \begin{align*}
      \frac{I_{OUT}}{I_{IN}} &= \frac{I_{M_1}}{I_{M_0}} = \frac{k_{n_1} (V_{X} - V_{ee} - V_{th})^2}{k_{n_0} (V_{X} - V_{ee} - V_{th})^2} \hspace{1cm} \text{(From (5), (6): same } V_{GS}\text{)} \\
      &= \frac{k_{n_1}}{k_{n_0}} = \frac{\frac{K_P}{2}\frac{W_1}{L_1}}{\frac{K_P}{2}\frac{W_0}{L_0}} = \frac{\frac{W_1}{L_1}}{\frac{W_0}{L_0}} = \frac{W_1 L_0}{L_1 W_0}
    \end{align*}
    \(M_2\) carries \(I_{M_1}\) and \(M_3\) carries \(I_{M_0}\) (from (2), (3)); choosing
    \(\frac{W_2 L_3}{L_2 W_3} = \frac{W_1 L_0}{L_1 W_0}\) gives \(M_2\) and \(M_3\) the same
    \(V_{ov}\) (from (7), (8)), so \(V_Y = V_Z - V_{th} - V_{ov_2} = V_X\): \(M_1\) then sees
    \(V_{DS_1} = V_{DS_0}\) and the ratio holds even with \(K_{CLM} \neq 0\).
  \item \textbf{Small Signal, Frequency Domain, AC Input Analysis (Q4)}:\\
    Applying KCL at $V_{out}$:
    \begin{align*}
      V_{out}G_5 + (V_{out} - V_{y})g_{o_2} + (V_{z} - V_{y})g_{m_2} &= 0 \\
      V_{out}G_5 + V_{out}g_{o_2} - V_{y}g_{o_2} + V_{z}g_{m_2} - V_{y}g_{m_2} &= 0 \\
      V_{out}(G_5 + g_{o_2}) - V_{y}(g_{o_2} + g_{m_2}) + V_{z}g_{m_2} &= 0 \tag{9}
    \end{align*}
    Applying KCL at $V_{x}$:
    \begin{align*}
      V_{x}g_{o_0} + I_{G_1} + V_{x}g_{m_0} + (V_{x} - V_{z})g_{o_3} - (V_{z} - V_{x})g_{m_3} &= 0 \\
      V_{x}g_{o_0} + V_{x}g_{m_0} + (V_{x} - V_{z})g_{o_3} - (V_{z} - V_{x})g_{m_3} &= 0 \quad \text{[as } I_{G} \approx 0 \text{ by MOS physics]} \\
      V_{x}g_{o_0} + V_{x}g_{m_0} + V_{x}g_{o_3} - V_{z}g_{o_3} - V_{z}g_{m_3} + V_{x}g_{m_3} &= 0 \\
      V_{x}(g_{o_0} + g_{m_0} + g_{o_3} + g_{m_3}) - V_{z}(g_{o_3} + g_{m_3}) &= 0 \tag{10}
    \end{align*}
    Applying KCL at $V_{y}$:
    \begin{align*}
      V_{y}g_{o_1} + V_{x}g_{m_1} + (V_{y} - V_{out})g_{o_2} - (V_{z} - V_{y})g_{m_2} &= 0 \\
      V_{y}g_{o_1} + V_{x}g_{m_1} + V_{y}g_{o_2} - V_{out}g_{o_2} - V_{z}g_{m_2} + V_{y}g_{m_2} &= 0 \\
      V_{y}(g_{o_1} + g_{o_2} + g_{m_2}) - V_{out}g_{o_2} - V_{z}g_{m_2} + V_{x}g_{m_1} &= 0 \tag{11}
    \end{align*}
    Applying KCL at $V_{z}$:
    \begin{align*}
      V_{z}G_4 + (V_{z} - V_{x})g_{o_3} + (V_{z} - V_{x})g_{m_3} + I_{G_2} &= 0 \\
      V_{z}G_4 + (V_{z} - V_{x})g_{o_3} + (V_{z} - V_{x})g_{m_3} &= 0 \quad \text{[as } I_{G} \approx 0 \text{ by MOS physics]} \\
      V_{z}G_4 + V_{z}g_{o_3} - V_{x}g_{o_3} + V_{z}g_{m_3} - V_{x}g_{m_3} &= 0 \\
      V_{z}(G_4 + g_{o_3} + g_{m_3}) - V_{x}(g_{o_3} + g_{m_3}) &= 0 \tag{12}
    \end{align*}
    \paragraph{Voltage Transfer \(\frac{V_{out}}{V_z}\):} a source \(V_s\) forcing \(V_z = V_s\)
    supplies whatever current (12) needs, so (12) drops out and (9), (10),
    (11) are solved with \(V_z\) known. Let
    \begin{align*}
      S &= g_{m_0} + g_{o_0} + g_{m_3} + g_{o_3} \hspace{1cm} P = g_{m_2} + g_{o_1} + g_{o_2} \hspace{1cm} a = \frac{g_{m_3} + g_{o_3}}{S}
    \end{align*}
    From (10):
    \begin{align*}
      V_x S &= V_z (g_{o_3} + g_{m_3}) \\
      V_x &= a V_z \tag{13}
    \end{align*}
    From (11) with (13):
    \begin{align*}
      V_y P &= V_{out} g_{o_2} + V_z g_{m_2} - a V_z g_{m_1} \\
      V_y &= \frac{V_{out} g_{o_2} + V_z (g_{m_2} - a g_{m_1})}{P} \tag{14}
    \end{align*}
    Substituting (14) into (9) and multiplying by \(P\):
    \begin{align*}
      V_{out}(G_5 + g_{o_2}) P - (g_{o_2} + g_{m_2})\left(V_{out} g_{o_2} + V_z (g_{m_2} - a g_{m_1})\right) + V_z g_{m_2} P &= 0 \\
      V_{out}\left[(G_5 + g_{o_2}) P - g_{o_2}(g_{o_2} + g_{m_2})\right] &= V_z\left[(g_{o_2} + g_{m_2})(g_{m_2} - a g_{m_1}) - g_{m_2} P\right]
    \end{align*}
    The left bracket, with \(P - (g_{o_2} + g_{m_2}) = g_{o_1}\):
    \begin{align*}
      G_5 P + g_{o_2} P - g_{o_2}(g_{o_2} + g_{m_2}) &= G_5 P + g_{o_2}\left(P - g_{o_2} - g_{m_2}\right) = G_5 P + g_{o_1} g_{o_2}
    \end{align*}
    The right bracket:
    \begin{align*}
      (g_{o_2} + g_{m_2}) g_{m_2} - a g_{m_1}(g_{o_2} + g_{m_2}) - g_{m_2} P
        &= -g_{m_2}\left(P - g_{o_2} - g_{m_2}\right) - a g_{m_1}(g_{m_2} + g_{o_2}) \\
        &= -g_{m_2} g_{o_1} - a g_{m_1}(g_{m_2} + g_{o_2})
    \end{align*}
    so, multiplying numerator and denominator by \(S\) (\(aS = g_{m_3} + g_{o_3}\)):
    \begin{align*}
      \frac{V_{out}}{V_z} &= -\frac{g_{m_2} g_{o_1} + a g_{m_1}(g_{m_2} + g_{o_2})}{G_5 P + g_{o_1} g_{o_2}} \\
      \frac{V_{out}}{V_z} &= -\frac{g_{m_1}(g_{m_2} + g_{o_2})(g_{m_3} + g_{o_3}) + g_{m_2} g_{o_1}(g_{m_0} + g_{o_0} + g_{m_3} + g_{o_3})}{(g_{m_0} + g_{o_0} + g_{m_3} + g_{o_3})\left(G_5 (g_{m_2} + g_{o_1} + g_{o_2}) + g_{o_1} g_{o_2}\right)}
    \end{align*}
    \paragraph{Current Transfer \(\frac{I_{out}}{I_{in}}\):} with the output shorted
    (\(V_{out} = 0\)), \(I_{in}\) is the current \(M_3\) delivers into \(V_x\) and \(I_{out}\) the
    current \(M_1\) draws from \(V_y\). KCL at \(V_x\) says all of \(I_{in}\) flows down
    \(M_0\), and KCL at \(V_y\) says all of \(I_{out}\) flows down \(M_1\):
    \begin{align*}
      I_{in} &= (g_{m_0} + g_{o_0}) V_x \hspace{1cm} I_{out} = g_{m_1} V_x + g_{o_1} V_y
    \end{align*}
    From (14) with \(V_{out} = 0\) and \(V_z = \frac{V_x}{a} = V_x\frac{S}{g_{m_3} + g_{o_3}}\) from (13):
    \begin{align*}
      V_y &= \frac{V_z g_{m_2} - V_x g_{m_1}}{P} = \frac{V_x}{P}\left(\frac{g_{m_2} S}{g_{m_3} + g_{o_3}} - g_{m_1}\right) \\
      I_{out} &= g_{m_1} V_x + \frac{g_{o_1} V_x}{P}\left(\frac{g_{m_2} S}{g_{m_3} + g_{o_3}} - g_{m_1}\right)
               = \frac{V_x}{P}\left(g_{m_1}(P - g_{o_1}) + \frac{g_{o_1} g_{m_2} S}{g_{m_3} + g_{o_3}}\right) \\
      I_{out} &= V_x\frac{g_{m_1}(g_{m_2} + g_{o_2})(g_{m_3} + g_{o_3}) + g_{m_2} g_{o_1} S}{P (g_{m_3} + g_{o_3})} \\
      \frac{I_{out}}{I_{in}} &= \frac{g_{m_1}(g_{m_2} + g_{o_2})(g_{m_3} + g_{o_3}) + g_{m_2} g_{o_1} S}{(g_{m_0} + g_{o_0})(g_{m_3} + g_{o_3})(g_{m_2} + g_{o_1} + g_{o_2})}
    \end{align*}
    With \(g_m \gg g_o\), keeping the largest term in each factor:
    \(\frac{I_{out}}{I_{in}} \approx \frac{g_{m_1} g_{m_2} g_{m_3}}{g_{m_0} g_{m_3} g_{m_2}} = \frac{g_{m_1}}{g_{m_0}}\), the
    Q5 ratio \(\frac{W_1 L_0}{L_1 W_0}\) for equal \(V_{ov}\).
\end{enumerate}


\subsubsection{Analysis}
\begin{enumerate}
  \item \textbf{Output Load (Max)}: \(M_2\) saturated needs
    \(V_{DS_2} \geq V_{ov_2}\). With \(V_{GS_2} = V_Z - V_Y = V_{th} + V_{ov_2}\), i.e.
    \(V_Y = V_Z - V_{th} - V_{ov_2}\):
    \begin{align*}
      V_{OUT,\min} &= V_Y + V_{ov_2} = V_Z - V_{th} - V_{ov_2} + V_{ov_2} = V_Z - V_{th} \\
      V_{cc} - I_{M_2} R_{L} &\geq V_{OUT,\min} \\
      R_{L, \max} &= \frac{V_{cc} - V_{OUT,\min}}{I_{M_2}}
    \end{align*}
  \item \textbf{Impedances}:
    \paragraph{At Input:} a test current \(I_{test}\) into \(V_z\) adds to the
    left side of (12):
      \begin{align*}
        V_{z}(G_4 + g_{o_3} + g_{m_3}) - V_{x}(g_{o_3} + g_{m_3}) &= I_{test} \tag{15}
      \end{align*}
      (10) still gives \(V_x = a V_z\) from (13), so
      \begin{align*}
        I_{test} &= V_z\left(G_4 + (g_{o_3} + g_{m_3}) - a(g_{o_3} + g_{m_3})\right) = V_z\left(G_4 + (g_{m_3} + g_{o_3})(1 - a)\right) \\
        1 - a &= \frac{S - (g_{m_3} + g_{o_3})}{S} = \frac{g_{m_0} + g_{o_0}}{S} \\
        I_{test} &= V_z\left(G_4 + \frac{(g_{m_3} + g_{o_3})(g_{m_0} + g_{o_0})}{S}\right) = V_z\frac{G_4 S + (g_{m_3} + g_{o_3})(g_{m_0} + g_{o_0})}{S} \\
        Z_{in} &= \frac{V_{z}}{I_{test}} = \frac{g_{m_0} + g_{m_3} + g_{o_0} + g_{o_3}}{G_{4}(g_{m_0} + g_{m_3} + g_{o_0} + g_{o_3}) + (g_{m_0} + g_{o_0})(g_{m_3} + g_{o_3})}
      \end{align*}
      i.e.\ \(R_4 \parallel \left(\frac{1}{g_{m_0} + g_{o_0}} + \frac{1}{g_{m_3} + g_{o_3}}\right)\): two diodes in series
      shunted by \(R_4\).
    \paragraph{At Output:} a test current \(I_{test}\) into \(V_{out}\) adds to
      the left side of (9):
      \begin{align*}
        V_{out}(G_5 + g_{o_2}) - V_{y}(g_{o_2} + g_{m_2}) + V_{z}g_{m_2} &= I_{test} \tag{16}
      \end{align*}
      (10) and (12) are now both homogeneous in \(V_x, V_z\); from (10)
      \(V_x = aV_z\), and (12) becomes
      \begin{align*}
        V_z\left(G_4 + (g_{m_3} + g_{o_3})(1 - a)\right) &= 0
      \end{align*}
      The bracket is positive, so \(V_z = 0\) and \(V_x = 0\). Then (11) gives
      \begin{align*}
        V_y P &= V_{out} g_{o_2} \hspace{1cm} \Rightarrow \hspace{1cm} V_y = \frac{g_{o_2}}{P} V_{out}
      \end{align*}
      and (16):
      \begin{align*}
        I_{test} &= V_{out}(G_5 + g_{o_2}) - \frac{g_{o_2}(g_{o_2} + g_{m_2})}{P} V_{out}
                  = V_{out}\frac{G_5 P + g_{o_2}(P - g_{o_2} - g_{m_2})}{P} = V_{out}\frac{G_5 P + g_{o_1} g_{o_2}}{P} \\
        Z_{out} &= \frac{V_{out}}{I_{test}} = \frac{g_{m_2} + g_{o_1} + g_{o_2}}{G_{5}(g_{m_2} + g_{o_1} + g_{o_2}) + g_{o_1} g_{o_2}}
      \end{align*}
\end{enumerate}


\subsubsection{Design}
Large Signal, Time Domain, DC Input Analysis (Q5) can be used to set the bias
currents and voltages in this circuit.
\begin{enumerate}
  \item Choose appropriate transistor dimensions
    ($W$ and $L$) for transistors $M_0$ and $M_1$
    to achieve the desired current gain.
    \begin{align*}
      \frac{I_{OUT}}{I_{IN}} &= \frac{W_1 L_0}{L_1 W_0} = \frac{W_2 L_3}{L_2 W_3}
    \end{align*}
  \item Choose a suitable overdrive voltage value
    ($V_{ov}$) for computing the bias voltages at
    nodes $V_{X}$, $V_{Y}$, and $V_{Z}$.
    \begin{align*}
      V_{ov_0} &= V_{GS_0} - V_{th} \\
      V_{ov_0} &= V_{X} - V_{ee} - V_{th} \\
      V_{X} &= V_{ov_0} + V_{ee} + V_{th} \\
      \\
      V_{ov_3} &= V_{GS_3} - V_{th} \\
      V_{ov_3} &= V_{Z} - V_{X} - V_{th} \\
      V_{Z} &= V_{ov_3} + V_{X} + V_{th} \\
      \\
      V_{ov_2} &= V_{GS_2} - V_{th} \\
      V_{ov_2} &= V_{Z} - V_{Y} - V_{th} \\
      V_{Y} &= V_{Z} - V_{ov_2} - V_{th}
    \end{align*}
  \item Using the chosen $V_{ov}$, compute $W_0$
    to set the required bias current $I_{IN}$.
    \begin{align*}
      I_{IN} &\approx I_{M_0} = \frac{K_P}{2} \frac{W_0}{L_0} V_{ov_0}^2 \\
             &\approx I_{M_3} = \frac{K_P}{2} \frac{W_3}{L_3} V_{ov_3}^2 \\
      \therefore W_0 &= \frac{2 I_{IN} L_0}{K_P V_{ov_0}^2} \\
                 W_3 &= \frac{2 I_{IN} L_3}{K_P V_{ov_3}^2}
    \end{align*}
    $W_1$ can then be computed using the current
    gain relation from step (a). Similarly, $W_2$
    can be computed using the current gain relation
    from step (a).
  \item Compute the required value of $G_4$ (or
    $R_4$) to ensure that the voltage at node
    $V_{Z}$ is as calculated in step (b).
    \begin{align*}
      V_{Z} &= V_{cc} - \frac{(I_{M_3} + I_{G_2})}{G_4} \\
      V_{Z} &\approx V_{cc} - \frac{I_{M_3}}{G_4} \quad \text{[assuming } I_{G_2} \approx 0 \text{ by MOS physics]} \\
      V_{Z} - V_{cc} &\approx - \frac{I_{M_3}}{G_4} \\
      \frac{(I_{M_3})}{(V_{cc} - V_{Z})} &\approx G_4 \\
      \frac{(I_{IN})}{(V_{cc} - V_{Z})} &\approx G_4 \quad \text{[} I_{IN} = I_{M_3} = I_{M_0} \text{]} \\
      \therefore R_4 &\approx \frac{V_{cc} - V_{Z}}{I_{IN}}
    \end{align*}
\end{enumerate}

\subsubsection{Question}
Design a cascode current mirror with a current
gain of 5. The input current is $I_{IN} = 10 \mu$A.
Assume the following specifications:
\begin{enumerate}
  \item $V_{cc} = +5$ V, $V_{ee} = 0$ V
  \item $V_{th} = 0.7$ V
  \item $K_P = 100 \mu A/V^2$
  \item $K_{CLM} = 0.02$ V$^{-1}$
  \item $L_{\min} = 1 \mu m$
  \item $V_{L, \max} = V_{cc}$
  \item $V_{ov} = 0.2$ V for all transistors
\end{enumerate}
Find the following:
\begin{enumerate}
  \item Transistor dimensions ($W$ and $L$) for all
    transistors.
  \item Values of resistor $R_4$ required to set
    the bias current.
  \item Output voltage range and maximum output
    load that can be connected at the output.
  \item Small signal output resistance of the
    current mirror. Assume a load resistance
    $R_L = 50 k\Omega$ is connected at the
    output.
\end{enumerate}
\paragraph{Solution:}
\begin{enumerate}
  \item Transistor dimensions:
    \begin{align*}
      \frac{I_{OUT}}{I_{IN}} &= \frac{W_1 L_0}{L_1 W_0} = \frac{W_2 L_3}{L_2 W_3} = 5 \\
      \\
      V_{X} &= V_{ov_0} + V_{ee} + V_{th} = 0.2 + 0 + 0.7 = 0.9 \text{ V} \\
      V_{Z} &= V_{ov_3} + V_{X} + V_{th} = 0.2 + 0.9 + 0.7 = 1.8 \text{ V} \\
      V_{Y} &= V_{Z} - V_{ov_2} - V_{th} = 1.8 - 0.2 - 0.7 = 0.9 \text{ V} \\
      \\
      W_0 &= \frac{2 I_{IN} L_0}{K_P V_{ov_0}^2} = \frac{2 \times 10 \times 10^{-6} \times 1 \times 10^{-6}}{100 \times 10^{-6} \times (0.2)^2} = 5 \mu m \\
      W_3 &= \frac{2 I_{IN} L_3}{K_P V_{ov_3}^2} = \frac{2 \times 10 \times 10^{-6} \times 1 \times 10^{-6}}{100 \times 10^{-6} \times (0.2)^2} = 5 \mu m \\
      W_1 &= 5 \times \frac{W_0}{L_0} \times L_1 = 5 \times \frac{5}{1} \times 1 = 25 \mu m \\
      W_2 &= 5 \times \frac{W_3}{L_3} \times L_2 = 5 \times \frac{5}{1} \times 1 = 25 \mu m
    \end{align*}
  \item Value of resistor $R_4$:
    \begin{align*}
      R_4 &\approx \frac{V_{cc} - V_{Z}}
                      {I_{IN}} = \frac{5 - 1.8}{10 \times 10^{-6}} = 320 k\Omega
    \end{align*}
  \item Output voltage range and maximum output load
    (\(I_{M_2} = I_{OUT} = 5 \cdot 10\,\mu A = 50\,\mu A\)):
    \begin{align*}
      V_{OUT,\min} &\approx V_{Y} + V_{ov_2} = 0.9 + 0.2 = 1.1 \text{ V} \\
      V_{OUT,\max} &\approx V_{cc} = 5 \text{ V} \\
      R_{L, \max} &= \frac{V_{cc} - V_{OUT,\min}}{I_{M_2}} = \frac{5 - 1.1}{50 \times 10^{-6}} = 78 k\Omega
    \end{align*}
    With \(R_L = 50\,k\Omega\): \(V_{OUT} = V_{cc} - I_{M_2} R_L = 5 - 50 \times 10^{-6} \times 50 \times 10^{3} = 5 - 2.5 = 2.5\,V \geq 1.1\,V\),
    so the \(50\,k\Omega\) load keeps \(M_2\) saturated.
  \item Small signal output resistance:
    \begin{align*}
      g_{m_2} &= \frac{2 I_{M_2}}{V_{ov_2}} = \frac{2 \times 50 \times 10^{-6}}{0.2} = 5 \times 10^{-4} \text{ S} \\
      g_{o_1} &= I_{M_1} K_{CLM} = 50 \times 10^{-6} \times 0.02 = 1 \times 10^{-6} \text{ S} \\
      g_{o_2} &= I_{M_2} K_{CLM} = 50 \times 10^{-6} \times 0.02 = 1 \times 10^{-6} \text{ S} \\
      G_5 &= \frac{1}{R_L} = \frac{1}{50 k\Omega} = 2 \times 10^{-5} \text{ S} \\
      \\
      g_{m_2} + g_{o_1} + g_{o_2} &= 5 \times 10^{-4} + 1 \times 10^{-6} + 1 \times 10^{-6} = 5.02 \times 10^{-4} \text{ S} \\
      Z_{out} &= \frac{g_{m_2} + g_{o_1} + g_{o_2}}{G_{5}(g_{m_2} + g_{o_1} + g_{o_2}) + g_{o_1} g_{o_2}} \\
              &= \frac{5.02 \times 10^{-4}}{(2 \times 10^{-5})(5.02 \times 10^{-4}) + (10^{-6})(10^{-6})} \\
              &= \frac{5.02 \times 10^{-4}}{1.004 \times 10^{-8} + 10^{-12}} = \frac{5.02 \times 10^{-4}}{1.0041 \times 10^{-8}} \approx 49.995 k\Omega
    \end{align*}
    The load dominates. The mirror alone (\(G_5 \to 0\)):
    \begin{align*}
      Z_{out,mirror} &= \frac{g_{m_2} + g_{o_1} + g_{o_2}}{g_{o_1} g_{o_2}} = \frac{5.02 \times 10^{-4}}{10^{-12}} = 502 M\Omega \\
      g_{m_2} r_{o_1} r_{o_2} &= \frac{g_{m_2}}{g_{o_1} g_{o_2}} = \frac{5 \times 10^{-4}}{10^{-12}} = 500 M\Omega \hspace{1cm} \text{(the \(g_m \gg g_o\) estimate)}
    \end{align*}
\end{enumerate}

\subsubsection{Code}
\begin{lstlisting}[language=Octave]
  clear; clc;
  pkg load symbolic;

  syms Vx Vy Vz Vout Itest
  syms G5 G4 go0 gm0 go1 gm1 go2 gm2 go3 gm3

  % small signal kcl, equations (9) to (12)
  eq9  = Vout*(G5 + go2) - Vy*(go2 + gm2) + Vz*gm2 == 0;  % KCL at Vout
  eq10 = Vx*(go0 + gm0 + go3 + gm3) - Vz*(go3 + gm3) == 0; % KCL at Vx
  eq11 = Vy*(go1 + go2 + gm2) - Vout*go2 - Vz*gm2 + Vx*gm1 == 0; % KCL at Vy
  eq12 = Vz*(G4 + go3 + gm3) - Vx*(go3 + gm3) == 0;        % KCL at Vz

  % voltage transfer: Vz forced by a source, so (12) drops out
  sol = solve([eq9, eq10, eq11], [Vout, Vx, Vy]);
  printf('=== Vout/Vz ===\n'); pretty(factor(simplify(sol.Vout/Vz)))

  % current transfer: output shorted (Vout = 0), Vz forced
  sol = solve([eq10, subs(eq11, Vout, 0)], [Vx, Vy]);
  Iin  = (gm0 + go0)*sol.Vx;            % all of Iin flows down M0
  Iout = gm1*sol.Vx + go1*sol.Vy;       % all of Iout flows down M1
  printf('=== Iout/Iin ===\n'); pretty(factor(simplify(Iout/Iin)))

  % input impedance: Itest injected into Vz, equation (15)
  eq15 = Vz*(G4 + go3 + gm3) - Vx*(go3 + gm3) == Itest;
  sol = solve([eq9, eq10, eq11, eq15], [Vout, Vx, Vy, Vz]);
  printf('=== Zin ===\n'); pretty(factor(simplify(sol.Vz/Itest)))

  % output impedance: Itest injected into Vout, equation (16)
  eq16 = Vout*(G5 + go2) - Vy*(go2 + gm2) + Vz*gm2 == Itest;
  sol = solve([eq16, eq10, eq11, eq12], [Vout, Vx, Vy, Vz]);
  printf('=== Zout ===\n'); pretty(factor(simplify(sol.Vout/Itest)))
\end{lstlisting}
